\addchap{Use of artificial intelligence-based tools}
\setcounter{chapter}{1}

%Im Rahmen dieser Arbeit wurden Künstliche Intelligenz (KI)\index{Künstliche Intelligenz} basierte Werkzeuge benutzt. Tabelle~\ref{tab:anhang_uebersicht_KI_werkzeuge} gibt eine Übersicht über die verwendeten Werkzeuge und den jeweiligen Einsatzzweck.
Artificial intelligence (AI)\index{Artificial intelligence} based tools were used in this work. Table~\ref{tab:anhang_uebersicht_KI_werkzeuge} provides an overview of the tools used and their respective purpose.

\begin{table}[hbt]	
	\centering
	\renewcommand{\arraystretch}{1.5}	% Scales the row height of the table
	\captionabove[List of AI-based tools used]{List of AI-based tools used}
	\label{tab:appendix_overview_AI_tools}
	\begin{tabular}{>{\raggedright\arraybackslash}p{0.3\linewidth} >{\raggedright\arraybackslash}p{0.65\linewidth}}
		\textbf{Tool} & \textbf{Description of Use}\\
		\hline 
		\hline
		Mercedes-Benz Direct Chat & Assistance with phrasing and spell checking (see entire thesis) \\
		DeepL & Translation of the thesis \\
		ChatGPT & Content structuring, summarization and technical formulation (see entire thesis) \\
	\end{tabular} 
\end{table}

%\addchap{B Ergänzungen}
%\setcounter{chapter}{2}

%\section{Details zu bestimmten theoretischen Grundlagen}

%\section{Weitere Details, welche im Hauptteil den Lesefluss behindern}

%\addchap{C Details zu Laboraufbauten und Messergebnissen}
%\setcounter{chapter}{3}
%\setcounter{section}{0}
%\setcounter{table}{0}
%\setcounter{figure}{0}

%\section{Versuchsanordnung}

%\section{Liste der verwendeten Messgeräte}

%\section{Übersicht der Messergebnisse}

%\section{Schaltplan und Bild der Prototypenplatine}

%\addchap{D Zusatzinformationen zu verwendeter Software}
%\setcounter{chapter}{4}
%\setcounter{section}{0}
%\setcounter{table}{0}
%\setcounter{figure}{0}

%\section{Struktogramm des Programmentwurfs}

%\section{Wichtige Teile des Quellcodes}

%\addchap{E Datenblätter}
%\setcounter{chapter}{5}
%\setcounter{section}{0}
%\setcounter{table}{0}
%\setcounter{figure}{0}

%\section{Einbinden von PDF-Seiten aus anderen Dokumenten}

%Auf den folgenden Seiten wird eine Möglichkeit gezeigt, wie aus einem anderen PDF-Dokument komplette Seiten übernommen werden können, z.~B. zum Einbindungen von Datenblättern. Der Nachteil dieser Methode besteht darin, dass sämtliche Formateinstellungen (Kopfzeilen, Seitenzahlen, Ränder, etc.) auf diesen Seiten nicht angezeigt werden. Die Methode wird deshalb eher selten gewählt. Immerhin sorgt das Package \textit{\glqq pdfpages\grqq}~für eine korrekte Seitenzahleinstellung auf den im Anschluss folgenden \glqq nativen\grqq~\LaTeX-Seiten.

%Eine bessere Alternative ist, einzelne Seiten mit \textit{\glqq$\backslash$includegraphics\grqq}~einzubinden.

%\includepdf[pages={2-4}]{docs/EingebundenesPDF.pdf}

%\clearpage
